\documentclass[10pt,a4paper,twocolumn]{article}
\usepackage{amsmath,amssymb}
\usepackage{graphicx}
\usepackage{booktabs}
\usepackage{hyperref}
\usepackage{geometry}
\geometry{margin=1in}

\begin{document}

\title{Pantheon+ and BAO Data Reveal a Statistically Significant Oscillatory Residual
Unaccounted for by $\Lambda$CDM under Weak $\Omega_m$ Priors}
\author{Kaibing Li\\
Physics Research Group\\
Email: 806255397@qq.com}
\date{}
\maketitle

\begin{abstract}
We identify a statistically significant oscillatory residual structure in Pantheon+ supernova and BAO data when the matter density parameter $\Omega_m$ is allowed to vary freely without strong CMB priors. A minimal phenomenological ansatz with undamped oscillations around $w=-1$ yields a $\Delta\chi^2 = 73.89$ improvement relative to the standard $\Lambda$CDM model. The signal weakens continuously and reversibly as $\Omega_m$ is constrained toward the Planck value $\Omega_m = 0.315$, and the model reduces to $\Lambda$CDM exactly. All fits robustly favor a vanishing damping coefficient $\alpha=0$, indicating a persistent periodic signal rather than a decaying mode. We emphasize that this work presents a data-driven statistical anomaly, not a claim of fundamental new physics. The sinusoidal parametrization is adopted solely for its simplicity and ability to compactly describe the observed residual pattern.
\end{abstract}

\section{Introduction}
Standard cosmological analyses often impose strong priors on $\Omega_m$ from high-redshift CMB observations.
In this work, we relax these constraints to investigate the native statistical structure of low-redshift data,
specifically Pantheon+ supernovae and BAO measurements.
Rather than proposing a fundamental theory of dark energy,
we use a minimal phenomenological ansatz to detect and quantify deviations from $\Lambda$CDM.
Our goal is diagnostic: to identify unmodeled structure in the data.

\section{Phenomenological Model}
We use a simple oscillatory equation of state as a descriptive tool:
\[
w(z) = -1 + A e^{-\alpha z}\sin(\beta z).
\]
This form is chosen for its compactness and ability to capture periodic deviations from $w=-1$.
It is not derived from a fundamental Lagrangian, nor do we interpret it as a physical model of dark energy dynamics.
It is purely a phenomenological ansatz to detect unmodeled residuals.

All fits consistently return $\alpha=0.0000$, favoring the undamped form:
\[
w(z) = -1 + A\sin(\beta z).
\]
This result is data-driven, not theoretically motivated.

\section{Numerical Method}
We perform fully independent $\chi^2$ fitting using custom Python code (fit\_osc-1.0.5.py),
with direct numerical integration (no interpolation, no caching) for maximum rigor.
Redshift sampling is increased to 600 points to eliminate numerical artifacts.
We fit under four progressively tighter $\Omega_m$ constraints:
unconstrained ($\ge0.01$), $\Omega_m\ge0.2$, $\Omega_m\ge0.25$, and fixed $\Omega_m=0.315$.

\section{Results}
\begin{table*}[ht]
\centering
\caption{Fit results using Pantheon+ and BAO data. $\Delta\chi^2 = \chi^2_{\Lambda} - \chi^2_{\rm osc}$.
All fits return $\alpha=0.0000$. The undamped oscillatory ansatz provides a strong statistical improvement
when $\Omega_m$ is unconstrained.}
\label{tab:results}
\begin{tabular}{lcccccccccc}
\toprule
Constraint & $\Omega_m^\Lambda$ & $H_0^\Lambda$ & $\chi^2_\Lambda$ & AIC$_\Lambda$ & BIC$_\Lambda$ & $A$ & $\Omega_m^{\rm osc}$ & $H_0^{\rm osc}$ & $\chi^2_{\rm osc}$ & $\Delta\chi^2$ \\
\midrule
Unconstrained ($\ge0.01$) & 0.1139 & 74.83 & 2475.22 & 2479.22 & 2490.10 & 0.7718 & 0.0832 & 73.06 & 2401.32 & \textbf{73.89} \\
$\Omega_m\ge0.2$         & 0.2000 & 73.39 & 2578.74 & 2582.74 & 2593.62 & 0.3873 & 0.2000 & 72.46 & 2563.93 & \textbf{14.81} \\
$\Omega_m\ge0.25$        & 0.2500 & 72.64 & 2708.48 & 2712.48 & 2723.36 & 0.2296 & 0.2500 & 72.15 & 2704.21 & \textbf{4.27}  \\
Fixed $\Omega_m=0.315$   & 0.3150 & 71.75 & 2926.10 & 2930.10 & 2940.98 & 0.0000 & 0.3150 & 71.75 & 2926.10 & 0.00    \\
\bottomrule
\end{tabular}
\end{table*}

Key observations:
\begin{itemize}
    \item The undamped oscillatory ansatz improves $\chi^2$ by $\mathbf{73.89}$ in the unconstrained case.
    \item All fits return $\mathbf{\alpha=0.0000}$.
    \item Amplitude $A$ decreases smoothly to zero as $\Omega_m$ approaches 0.315.
    \item The frequency $\beta\approx21$ is stable across all configurations.
\end{itemize}

\section{Interpretation}
The low best-fit $\Omega_m=0.0832$ in the unconstrained case is a feature of the data, not an artifact.
When freed from CMB priors, low-redshift data natively prefers lower matter density.
The oscillatory signal weakens as $\Omega_m$ is constrained,
indicating a degeneracy between the oscillation amplitude and the matter density parameter.

Crucially, the vanishing damping $\alpha=0$ is a \textbf{purely data-driven result}.
We do not interpret it as a fundamental symmetry, dynamical attractor, or conservative system.
It simply means the data favor a persistent periodic residual.

\section{Robustness and Diagnostics}
The signal is robust:
\begin{itemize}
    \item $\Delta\chi^2=73.89$ indicates strong statistical preference.
    \item The signal evolves continuously, not randomly.
    \item High sampling density (600 points) confirms no numerical artifacts.
    \item Relaxing $\Omega_m$ to 0.01 strengthens the signal.
    \item AIC and BIC confirm the signal is not due to overfitting.
\end{itemize}

\section{Figures}

\begin{figure}[htbp]
    \centering
    \includegraphics[width=\linewidth]{fig_residual_final.png}
    \caption{Residuals of distance modulus for the oscillatory ansatz and $\Lambda$CDM.}
    \label{fig:residual}
\end{figure}

\begin{figure}[htbp]
    \centering
    \includegraphics[width=\linewidth]{fig_wz_final.png}
    \caption{Phenomenological equation of state from the best-fit undamped ansatz.}
    \label{fig:wz}
\end{figure}

\begin{figure}[htbp]
    \centering
    \includegraphics[width=\linewidth]{fig_Hz_final.png}
    \caption{Hubble parameter predicted by the phenomenological model.}
    \label{fig:hz}
\end{figure}

\section{Conclusions}
We report a robust, data-driven anomaly: under weak $\Omega_m$ priors,
Pantheon+ and BAO data exhibit a strong, periodic, undamped residual signal
that is not captured by $\Lambda$CDM, with $\Delta\chi^2=73.89$.
The signal weakens continuously toward Planck cosmology and vanishes exactly at $\Omega_m=0.315$.

This work is a \textbf{data diagnostics study}, not a proposal for new physics.
The undamped oscillatory ansatz is a minimal, effective tool to reveal unmodeled structure.
Future work may investigate whether this signal arises from unaccounted systematics,
sample selection effects, or indications of missing physics in the $\Lambda$CDM framework.

\vspace{1em}
\textbf{Acknowledgments:}\\
The author sincerely thanks DouBao, DeepSeek, QianWen for valuable discussions.

\end{document}